\chapter{Antecedentes}
\label{cap:antecedentes}

\section{El cáncer como problema de salud pública y su impacto}

\subsection{Definición de cáncer}

El término cáncer, también conocido como tumor maligno o neoplasia maligna, no se refiere a una única enfermedad, sino que engloba un grupo amplio y heterogéneo de patologías que pueden originarse en prácticamente cualquier órgano o tejido del cuerpo humano (\cite{oms2022_cancer}).

Según la Organización Mundial de la Salud (OMS), el cáncer se caracteriza por la multiplicación rápida de células anormales que se extienden más allá de sus límites habituales y pueden invadir partes adyacentes del cuerpo o propagarse a otros órganos en un proceso denominado metástasis. En específico, esta enfermedad se produce cuando las células normales se transforman en células tumorales a través de un proceso de múltiples etapas, que generalmente consiste en la progresión desde una lesión precancerosa hacia un tumor maligno. Estas alteraciones son provocadas por la interacción entre factores genéticos propios de la persona afectada y agentes carcinógenos físicos (como radiaciones ultravioletas e ionizantes), químicos (como el tabaco, el asbesto, las aflatoxinas y el arsénico) y biológicos (como determinados virus, bacterias y parásitos) (\cite{oms2022_cancer}).

% oms2022_cancer: https://www.who.int/es/news-room/fact-sheets/detail/cancer

\subsection{Impacto global del cáncer}

El cáncer se ha convertido en una de las principales causas de mortalidad a nivel mundial. Según el Reporte Mundial de Cáncer 2020 de la OMS, esta enfermedad corresponde a la segunda causa más frecuente de mortalidad global (16,3\%), siendo superada únicamente por las enfermedades cardiovasculares (32\%). En particular, durante el año 2020 el cáncer constituyó la primera y segunda causa (en algunos países) de mortalidad prematura en 134 de 183 países. En América (incluyendo América del Norte, América del Sur y el Caribe), el cáncer se posicionó como la segunda causa principal de muerte, registrándose más de 1,4 millones de fallecimientos asociados a esta enfermedad, equivalentes al 14\% de las muertes oncológicas a nivel mundial. De estas defunciones, aproximadamente 713.414 correspondieron específicamente a Sudamérica y el Caribe (5\% del total mundial), destacando Brasil, México y Argentina como los países con mayor cantidad de casos (\cite{minsal2022}).

Respecto de los tipos de cáncer, aquellos que registraron el mayor número de muertes a nivel mundial según GLOBOCAN al año 2020 fueron el cáncer de pulmón (18\%), el cáncer colorrectal (9,4\%), el cáncer de hígado (8,3\%), el cáncer de estómago (7,7\%) y el cáncer de mama (6,9\%). En relación con la mortalidad oncológica según sexo, en mujeres la mayor tasa de defunción corresponde al cáncer de mama, con una tasa de mortalidad de 13,6 por cada 100.000 mujeres, seguido del cáncer de pulmón (11,2), el cáncer cervicouterino (7,3) y el cáncer colorrectal (7,2 por cada 100.000 mujeres). En hombres, predomina el cáncer de pulmón con una tasa de 25,9 por cada 100.000 hombres, seguido por el cáncer de hígado (12,9), el cáncer colorrectal y el cáncer de estómago, ambos con una tasa cercana a 11 por cada 100.000 hombres (\cite{minsal2022}).

En síntesis, la OMS estima que el cáncer ocasiona alrededor de 9 millones de muertes anuales a nivel mundial, representando aproximadamente una de cada seis muertes en el mundo. Además, cerca del 30\% de las muertes por cáncer se atribuyen a factores de riesgo modificables, como el consumo de tabaco y alcohol, el índice de masa corporal, la alimentación inadecuada y el sedentarismo, lo que evidencia un importante potencial de prevención mediante estrategias de salud pública orientadas a reducir su impacto (\cite{minsal2022}).

Siguiendo esta línea, a partir del año 2013 la OMS impulsó un plan de acción mundial para la prevención y control de enfermedades no transmisibles, estableciendo como meta reducir en un 25\% la mortalidad prematura causada por las cuatro principales enfermedades crónicas hacia el año 2025, entre ellas el cáncer. Este plan contempla estrategias centradas en prevención, diagnóstico precoz, acceso oportuno a tratamiento de calidad y monitoreo continuo de los resultados sanitarios (\cite{minsal2022}).

Respecto de las proyecciones epidemiológicas, se estima que, si las condiciones estratégicas y epidemiológicas se mantienen sin cambios significativos en el tiempo (asumiendo estabilidad demográfica, poblacional y en las tasas de incidencia), el número anual de nuevos casos diagnosticados en Sudamérica y el Caribe aumentará desde 1,4 millones de casos en 2018 a aproximadamente 2,5 millones en 2040, lo que representaría un incremento cercano al 78\% (\cite{minsal2022}).

\subsection{Epidemiología e impacto socioeconómico del cáncer en Chile}

En el transcurso de las últimas décadas, Chile ha experimentado una profunda transición demográfica y epidemiológica, caracterizada por el envejecimiento progresivo de la población y el aumento de las enfermedades crónicas no transmisibles (ECNT). En este contexto, el cáncer ha desplazado a las enfermedades cardiovasculares en múltiples regiones del país, consolidándose como una de las principales causas de mortalidad y morbilidad a nivel nacional (\cite{minsal2022}).


Respecto de los casos de cáncer registrados durante el año 2020, según GLOBOCAN (\textit{Global Cancer Observatory}), en mujeres los cánceres con mayor incidencia en Chile fueron el cáncer de mama, colorrectal, cervicouterino, de pulmón y de tiroides. Mientras que en hombres, los de mayor incidencia correspondieron al cáncer de próstata, colorrectal, de estómago, de pulmón y de riñón. En relación con la mortalidad oncológica, en mujeres los principales tipos de cáncer fueron mama, pulmón, colorrectal, estómago y cervicouterino, mientras que en hombres correspondieron al cáncer de estómago, pulmón, próstata, colorrectal e hígado (\cite{minsal2022}).

% minsal2022: https://www.iccp-portal.org/sites/default/files/plans/Plan-Nacional-de-Cancer-2022-2027.pdf

Por otro lado, según estimaciones recientes de la Organización para la Cooperación y el Desarrollo Económicos (OCDE), publicadas en su informe de 2024 sobre el impacto del cáncer en Chile, se proyecta que entre los años 2023 y 2050 el cáncer será responsable del 18\% de las muertes prematuras (antes de los 75 años) en el país, lo que equivale a que una de cada seis muertes prematuras será consecuencia directa de esta enfermedad. Esto se traduce en un estimado de 9.900 muertes prematuras anuales y una reducción promedio de 1,6 años en la esperanza de vida de la población chilena (\cite{oecd2024}).

Más allá de la carga de mortalidad y morbilidad que impone esta enfermedad, la atención oncológica ejerce una presión significativa sobre la infraestructura de la red asistencial pública y el presupuesto estatal. El mismo informe de la OCDE estima que el gasto per cápita en salud asociado al cáncer en Chile experimentará un incremento del 123\% hacia el año 2050, impulsado principalmente por el envejecimiento poblacional (\cite{oecd2024}). Frente a este escenario, el Ministerio de Salud (MINSAL) formuló el Plan Nacional del Cáncer 2022--2027, marco estratégico y normativo derivado de la Ley Nacional del Cáncer (Ley 21.258), el cual reconoce oficialmente las deficiencias históricas en el acceso a atención oncológica, la saturación de centros de alta complejidad y la necesidad de fortalecer los sistemas de registro e inteligencia sanitaria. Asimismo, este plan destaca la importancia de optimizar los recursos hospitalarios y la vigilancia epidemiológica con el objetivo de mitigar el impacto socioeconómico de la enfermedad y garantizar la sostenibilidad de la red pública de salud (\cite{minsal2022}).

% oecd2024: https://www.oecd.org/en/publications/tackling-the-impact-of-cancer-on-health-the-economy-and-society_40335421-en/chile_7e89d292-en.html#:~:text=At%20least%20three%20factors%20will%20drive%20up,and%202050%2C%20all%20other%20things%20being%20equal.


\subsection{Clasificación y caracterización de las neoplasias malignas}

Para abordar estadísticamente el impacto del cáncer, resulta fundamental comprender tanto su naturaleza biológica como su estandarización clínica. El cáncer no corresponde a una patología única, sino a un conjunto complejo y heterogéneo de más de cien enfermedades distintas, caracterizadas por una proliferación celular anormal, descontrolada e invasiva. Debido a esta heterogeneidad, dichas patologías fueron estandarizadas a nivel internacional mediante la Clasificación Internacional de Enfermedades en su décima revisión (CIE-10), elaborada por la OMS, donde las neoplasias ocupan el Capítulo II (códigos C00 a D48). Esta clasificación proporciona un lenguaje codificado universal que permite la interoperabilidad entre registros clínicos y administrativos en contextos hospitalarios (\cite{paho}).

% paho: https://ais.paho.org/classifications/chapters/pdf/volume1.pdf#:~:text=Los%20trabajos%20de%20la%20D%C3%A9cima%20Revisi%C3%B3n%20de,OMS%20para%20la%20Clasificaci%C3%B3n%20de%20En%2D%20fermedades.

Con base en esta codificación, las neoplasias malignas pueden dividirse en dos grandes categorías basales:

\begin{itemize}
	\item \textbf{Neoplasias sólidas (C00--C80):} corresponden a tumores originados en órganos o tejidos específicos, como el cáncer de pulmón, mama o próstata. Estas neoplasias constituyen la mayoría de los casos oncológicos y su manejo suele requerir intervenciones quirúrgicas complejas, radioterapia y estancias hospitalarias prolongadas (\cite{ministerio_salud}). 
	\item \textbf{Neoplasias del tejido linfático y hematopoyético (C81--C96):} incluyen patologías como leucemias, linfomas y mielomas, originadas en la médula ósea o en los ganglios linfáticos. Debido a su naturaleza sistémica, presentan patrones de consumo de recursos y requerimientos de infraestructura distintos a los de las neoplasias sólidas, incluyendo necesidades como aislamiento clínico o trasplantes de médula ósea. Además, pueden originarse simultáneamente en múltiples localizaciones (\cite{ministerio_salud}).
\end{itemize}

No obstante, el sistema CIE-10 contempla distintos bloques principales de neoplasias malignas que determinan la naturaleza clínica de cada episodio hospitalario (\cite{paho}):

\begin{itemize}
	\label{item:tipos_cancer}
    \item \textbf{C00-C14 (Neoplasias malignas de labio, cavidad bucal y faringe):} Incluyen tumores originados en el tracto aerodigestivo superior, los cuales son altamente asociados a factores de riesgo como el consumo de tabaco y alcohol.
    \item \textbf{C15-C26 (Neoplasias malignas de los órganos digestivos):} Agrupan tumores del esófago, estómago, colon, recto e hígado, los cuales representan una alta carga de morbilidad y suelen implicar resecciones quirúrgicas mayores.
    \item \textbf{C30-C39 (Neoplasias malignas de los órganos respiratorios e intratorácicos):} Principalmente el cáncer de pulmón y laringe, caracterizados por una alta letalidad y requerimientos de soporte ventilatorio o cuidados críticos.
    \item \textbf{C40-C41 (Neoplasias malignas de los huesos y de los cartílagos articulares):} Tumores óseos primarios, como el osteosarcoma, que afectan la movilidad funcional del paciente y requieren cirugías traumatológicas complejas.
    \item \textbf{C43-C44 (Melanoma y otras neoplasias malignas de la piel):} Tumores cutáneos de diverso grado de agresividad, en donde el melanoma destaca por su alto potencial metastásico sistémico.
    \item \textbf{C45-C49 (Neoplasias malignas de los tejidos mesoteliales y de los tejidos blandos):} Incluyen mesoteliomas y diversos tipos de sarcomas originados en músculos, grasa o vasos sanguíneos.
    \item \textbf{C50 (Neoplasias malignas de la mama):} Una de las patologías oncológicas con mayor incidencia global en mujeres, cuyo abordaje es multidisciplinario (mastectomía, quimioterapia, hormonoterapia).
    \item \textbf{C51-C58 (Neoplasias malignas de los órganos genitales femeninos):} Tumores de cuello uterino, ovario y endometrio, fuertemente vinculados a planes de salud pública preventiva (como el Papanicolaou).
    \item \textbf{C60-C63 (Neoplasias malignas de los órganos genitales masculinos):} Dominado epidemiológicamente por el cáncer de próstata, cuyo manejo varía desde la vigilancia activa hasta la prostatectomía radical.
    \item \textbf{C64-C68 (Neoplasias malignas de las vías urinarias):} Neoplasias renales y de vejiga, que a menudo son diagnosticadas por hematuria y exigen procedimientos urológicos especializados.
    \item \textbf{C69-C72 (Neoplasias malignas del ojo, del encéfalo y de otras partes del sistema nervioso central):} Tumores cerebrales que, por su ubicación, presentan un riesgo neurológico crítico y demandan intervenciones neuroquirúrgicas de alta complejidad.
    \item \textbf{C73-C75 (Neoplasias malignas de la glándula tiroides y de otras glándulas endocrinas):} Tumores que afectan el sistema hormonal, generalmente con tasas de supervivencia altas si se detectan a tiempo.
    \item \textbf{C76-C80 (Neoplasias malignas de sitios mal definidos, secundarios y de sitios no especificados):} Códigos utilizados cuando el tumor primario es desconocido o se trata de una diseminación metastásica generalizada sin origen claro, lo que usualmente indica estados terminales o de diagnóstico tardío.
    \item \textbf{C81-C96 (Neoplasias malignas del tejido linfático, de los órganos hematopoyéticos y de tejidos afines):} Patologías sistémicas como leucemias y linfomas. Al no ser sólidas, su abordaje es predominantemente hematológico y farmacológico.
    \item \textbf{C97 (Neoplasias malignas de sitios múltiples independientes):} Casos clínicos excepcionales donde el paciente presenta dos o más tumores primarios malignos distintos desarrollados de forma simultánea.
\end{itemize}

% ministerio_salud: https://repositoriodeis.minsal.cl/ContenidoSitioWeb2020/uploads/2018/07/Orientaciones-de-codificaci%C3%B3n-para-el-Cap%C3%ADtulo-II-de-la-CIE-10.pdf

Respecto de la localización tumoral, el sistema CIE-10 introduce una distinción topográfica relevante entre localización primaria y localización secundaria (metástasis). Un tumor maligno primario corresponde al sitio original donde se inicia el crecimiento neoplásico, mientras que un tumor secundario representa la diseminación de células malignas hacia otros órganos o tejidos, lo que implica un estadio más avanzado de la enfermedad (\cite{paho}).

Esta diferenciación clínica constituye la base sobre la cual los sistemas de clasificación hospitalaria, como el estándar GRD, asignan en cada episodio hospitalario un diagnóstico principal (motivo central de ingreso) y diagnósticos secundarios (comorbilidades o metástasis), los cuales interactúan directamente para determinar la severidad del episodio, el riesgo de mortalidad y el nivel de intervenciones terapéuticas o quirúrgicas requeridas (\cite{ministerio_salud}).

%%%%%%%%%%%% -------------- %%%%%%%%%%%%%%%
\section{Sistema de Grupos Relacionados por el Diagnóstico (GRD)}

\subsection{Origen y evolución del modelo GRD en Chile}

El sistema de Grupos Relacionados por el Diagnóstico (GRD) fue desarrollado originalmente por investigadores de la Universidad de Yale durante la década de 1970, con el objetivo de clasificar los episodios de hospitalización en categorías clínicamente coherentes y con un consumo de recursos homogéneo, relacionando así la complejidad clínica de un hospital con su utilización de recursos. Su adopción a nivel internacional transformó la gestión sanitaria, permitiendo a los hospitales transitar desde modelos presupuestarios basados en promedios históricos hacia esquemas de financiamiento sustentados en la complejidad real de los casos atendidos. En síntesis, los GRD constituyen un sistema de clasificación de hospitalizaciones que permite analizar y caracterizar los distintos tipos de pacientes atendidos (\cite{Cid2024}).

% cid2024: https://iris.paho.org/server/api/core/bitstreams/87862a09-e2d2-48d3-b2a0-a42f7c160b84/content

En Chile, la implementación del sistema GRD comenzó mediante una serie de pilotos aislados orientados principalmente a la medición de la eficiencia clínica. En particular, durante el año 2001 los GRD fueron introducidos en un contexto académico; posteriormente, el Fondo Nacional de Salud (FONASA) comenzó a utilizarlos en el sector privado y desarrolló un piloto en el sector público en 2015. Tras casi dos décadas de avances y ajustes institucionales, en el año 2020 FONASA instauró formalmente el programa GRD como mecanismo oficial de pago para la atención cerrada en los establecimientos de la red pública (\cite{Cid2024}).

Posteriormente, durante el año 2023, mediante la Ley de Presupuestos del Sector Público, FONASA consolidó y amplió el financiamiento basado en GRD a 68 hospitales de alta complejidad, incorporando tres nuevos hospitales al programa. Estos cambios obligaron a los centros de salud a registrar con mayor precisión sus diagnósticos y procedimientos, con el fin de asegurar una correcta asignación de los recursos estatales (\cite{fonasa2023_presupuesto, metamodelo2024_analisis}).

% fonasa2023_presupuesto: https://www.senado.cl/site/presupuesto/2023/cumplimiento/Glosas%202023/16%20Salud/1602_19576_08112023_FONASA.pdf

% metamodelo2024_analisis: https://www.metamodelo.cl/2024/12/16/analisis-produccion-salud-2024/

\subsection{Estructura de datos y componentes del estándar IR-GRD}

Con el propósito de garantizar comparabilidad internacional y precisión clínica, Chile adoptó el estándar internacional conocido como International Refined Diagnosis Related Groups (IR-GRD), desarrollado por 3M Health Information Systems. Este modelo se basa en la clasificación de episodios hospitalarios en grupos clínicamente coherentes y con consumo homogéneo de recursos, utilizando variables demográficas, clínicas y hospitalarias para asignar cada episodio a un grupo específico. En particular, el principal insumo del algoritmo agrupador IR-GRD corresponde al Conjunto Mínimo Básico de Datos (CMBD) de hospitalización, definido como un registro estructurado que resume la información clínica y administrativa esencial del paciente al momento del egreso hospitalario (\cite{minsal2018_cmbd}).
% minsal2018_cmbd

% \cite{minsal2018_cmbd}: https://www.sigesa.com/grd/

Técnicamente, el CMBD se compone de múltiples dimensiones de variables, entre las cuales destacan:

\begin{itemize}
    \item \textbf{Variables demográficas y administrativas:} incluyen la edad, sexo, previsión de salud y comuna de residencia del paciente.
    \item \textbf{Variables clínicas (CIE-10):} comprenden el diagnóstico principal (motivo principal del ingreso) y los diagnósticos secundarios (comorbilidades, metástasis o complicaciones), codificados según el sistema CIE-10.
    \item \textbf{Indicadores de eficiencia:} incluyen los días de estadía (tiempo de hospitalización) y el peso relativo del GRD que corresponde a un valor numérico que refleja el costo esperado del episodio en comparación con el caso hospitalario promedio.
\end{itemize}

A partir de estas variables de entrada, el software agrupador procesa los episodios hospitalarios y genera dos índices categóricos fundamentales que operan como variables de salida u objetivo:


\begin{itemize}
	\item \textbf{Nivel de severidad o gravedad (GRS-S):} representa el nivel de descompensación o pérdida funcional de un órgano o sistema, considerando las características del paciente, los diagnósticos secundarios asociados al episodio y los procedimientos realizados. Este índice se encuentra estratificado en cuatro niveles: menor (1), moderado (2), mayor (3) y extremo (4) (\cite{revista_socampar}).
	\item \textbf{Riesgo de mortalidad (GRD-M):} corresponde a un indicador del riesgo de fallecimiento del paciente durante el episodio hospitalario, entendido como un índice descriptivo y no como un predictor clínico directo. Este índice se clasifica en los mismos cuatro niveles estadísticos utilizados para la severidad (\cite{revista_socampar}).
\end{itemize}

% \cite{revista_socampar}: https://revista-socampar.com/images/articulos/ArticulosSeparadosNum18/Original%203.pdf

\subsection{Relevancia de los procedimientos médicos y quirúrgicos en la gestión hospitalaria}

Además de los diagnósticos clínicos, el sistema GRD incorpora la codificación de los procedimientos médicos e intervenciones terapéuticas, diagnósticas o quirúrgicas realizadas durante la hospitalización del paciente, utilizando el sistema CIE-9-MC (Clasificación Internacional de Enfermedades, 9ª revisión, Modificación Clínica) (\cite{minsal2024_procedimientos}).
% (\cite{minsal2024_procedimientos}: https://www.salud.gob.sv/wp-content/uploads/2024/11/clasificacion-internacional-de-procedimientos-CIE-9MC-version-2024-lista-tabular-el-salvador-2024.pdf#:~:text=Como%20una%20instrucci%C3%B3n%20para%20codificar%20cada%20uno,es%C3%B3fago.%20Codificar%20adem%C3%A1s%20cualquier%20procedimiento%20simult%C3%A1neo%20de.

Respecto a la relevancia de los procedimientos, estos son fundamentales para capturar adecuadamente la complejidad clínica de los episodios hospitalarios. Por ejemplo, en pacientes oncológicos, dependiendo de la tipología y estadiaje del cáncer, pueden requerirse intervenciones específicas como resecciones tumorales, biopsias estereotáxicas, procedimientos reconstructivos o terapias invasivas. En consecuencia, los procedimientos reflejan las acciones terapéuticas concretas que determinan directamente mayores niveles de severidad, estancias hospitalarias prolongadas y un elevado consumo de recursos, dimensiones que precisamente busca gestionar y caracterizar el sistema GRD.

%%%%%%%%%%%%%%%%%%%%%%%%%%%%%%%%%%%%%%%%%%%%%%


\section{Marco teórico metodológico y computacional}

% \subsubsection{Selección de características y análisis de asociación estadística}

% COMPLETAR DESPUES

\subsection{Remuestreo estadístico y robustez mediante Bootstrapping}

El desbalance masivo de clases corresponde a un problema frecuente en bases de datos poblacionales de salud pública, donde la proporción de episodios que presentan una determinada patología (como cáncer) frente a aquellos que no la presentan es significativamente asimétrica. Para abordar un desafío de esta naturaleza sin incurrir en sesgos deterministas ni perder variabilidad clínica, existen técnicas de remuestreo computacional no paramétrico, como el \textit{bootstrapping}, introducido por Bradley Efron en 1979 (\cite{efron1979bootstrap}).

La técnica de remuestreo mediante \textit{bootstrapping} consiste en extraer aleatoriamente $B$ muestras repetidas de tamaño $N$, permitiendo el reemplazo de las observaciones; es decir, una misma observación puede ser seleccionada múltiples veces en distintas muestras. Matemáticamente, si se dispone de un conjunto original de datos $Z = (z_1, z_2, ..., z_n)$, una muestra \textit{bootstrap} $Z^{*b}$ contendrá elementos seleccionados desde $Z$ con probabilidad $1/n$ (\cite{efron1979bootstrap, hastie2009elements}).

En el contexto específico de esta investigación, el remuestreo se utiliza para forzar artificialmente el balance de clases, por ejemplo, emparejando aleatoriamente una determinada cantidad de pacientes oncológicos con pacientes sin cáncer en cada iteración. Al repetir este proceso sobre cientos de submuestras aleatorias, el algoritmo reduce el sesgo inherente asociado a una única partición de los datos. La justificación matemática de este enfoque metodológico reside en el Teorema del Límite Central (TLC) y en la Ley de los Grandes Números, los cuales establecen que, a medida que el número de iteraciones \textit{bootstrap} $B$ tiende al infinito, la distribución empírica de la importancia de las características converge hacia la verdadera distribución poblacional subyacente. Esto asegura que las variables finalmente seleccionadas correspondan a aquellas que demuestran persistencia estadística estructural, aislando el ruido estocástico y garantizando que los perfiles operativos identificados mediante selección de características no sean un artefacto del azar presente en una muestra aislada (\cite{hastie2009elements}).


\subsection{Modelos de aprendizaje automático y algoritmos de ensamble}

Para clasificar la información y extraer relaciones multivariadas entre las características de los datos GRD, este estudio emplea tres arquitecturas algorítmicas fundamentales: un modelo paramétrico utilizado como línea base (Regresión Logística) y dos algoritmos avanzados de ensamble basados en árboles de decisión, específicamente \textit{Random Forest} y XGBoost (\textit{eXtreme Gradient Boosting}) (\cite{hastie2009elements}).

% hastie2009elements: https://books.google.cl/books?id=VRzITwgNV2UC&pg=PA43&hl=es&source=gbs_selected_pages&cad=1#v=onepage&q&f=false

La \textbf{regresión logística} se seleccionó como modelo estadístico base debido a su elevada interpretabilidad inherente. Este algoritmo estima la probabilidad de que una instancia hospitalaria pertenezca a una determinada clase objetivo, como la condición de fallecido en el \textit{target} de mortalidad ($Y=1$), dada una matriz de características $X$, utilizando la función sigmoide definida en la Ecuación \ref{eq:logistica}:

\begin{equation}
    P(Y=1|X) = \frac{1}{1 + e^{-(\beta_0 + \beta_1 x_1 + \dots + \beta_n x_n)}}
    \label{eq:logistica}
\end{equation}

Donde $\beta_0$ corresponde al intercepto y $\beta_i$ representa los coeficientes que cuantifican la influencia de cada variable $x_i$ sobre los \textit{log-odds} de la variable objetivo (por ejemplo, mortalidad). El ajuste óptimo de estos coeficientes se obtiene mediante la minimización de la función de pérdida de entropía cruzada binaria o \textit{Log-Loss}, expresada en la Ecuación \ref{eq:logloss}, donde $y_k$ representa la etiqueta real del episodio hospitalario $k$ y $\hat{y}_k = P(Y=1|X_k)$ corresponde a la probabilidad estimada por el modelo para dicha instancia:

\begin{equation}
    \mathcal{L}(\beta) = -\frac{1}{N} \sum_{k=1}^{N} \left[ y_k \ln(\hat{y}_k) + (1 - y_k) \ln(1 - \hat{y}_k) \right]
    \label{eq:logloss}
\end{equation}

Para operar de manera estable en entornos de alta dimensionalidad y controlar la multicolinealidad estructural característica de las bases de datos de salud, la función de pérdida base se optimiza mediante técnicas de regularización elástica (\textit{Elastic Net}). Esta aproximación penaliza la magnitud de los coeficientes combinando simultáneamente regularización L1 (\textit{Lasso}) y L2 (\textit{Ridge}), transformando la función objetivo según la Ecuación \ref{eq:regularizacion}:

\begin{equation}
    \mathcal{L}_{\text{reg}}(\beta) = \mathcal{L}(\beta) + \lambda \left[ \alpha \sum_{j=1}^{m} |\beta_j| + \frac{1-\alpha}{2} \sum_{j=1}^{m} \beta_j^2 \right]
    \label{eq:regularizacion}
\end{equation}

En esta fórmula, el hiperparámetro $\lambda$ controla la intensidad global de la penalización aplicada al modelo, mientras que el parámetro de mezcla $\alpha \in [0,1]$ regula la proporción relativa entre ambos tipos de regularización. Estos términos poseen distintas implicancias operativas sobre el comportamiento de los coeficientes (\cite{hastie2009elements}):

\begin{itemize}
\item \textbf{Penalización L1 o Lasso ($\alpha = 1$):} añade un costo proporcional al valor absoluto de los coeficientes. Debido a la geometría de sus contornos normativos en forma de diamante, este operador induce dispersión (\textit{sparsity}), forzando que los coeficientes de variables redundantes o con baja asociación estadística se reduzcan exactamente a cero. En la práctica, este comportamiento permite que el algoritmo actúe automáticamente como un selector de características descriptivas.
\item \textbf{Penalización L2 o Ridge ($\alpha = 0$):} añade un costo proporcional al cuadrado de la magnitud de los coeficientes. Sus contornos elípticos contraen uniformemente los pesos hacia cero, pero sin anularlos completamente, como ocurre en la regularización L1. Esta penalización resulta especialmente efectiva para manejar grupos de variables correlacionadas, como diagnósticos concurrentes codificados, distribuyendo el peso de manera más equilibrada entre ellas.
\end{itemize}



\textbf{Random Forest} corresponde a un algoritmo de ensamble robusto basado en la técnica de agregación por arranque (\textit{Bagging} o \textit{Bootstrap Aggregating}), donde se construye un conjunto de $B$ árboles de decisión profundos de manera independiente ${T_1(X), T_2(X), \dots, T_B(X)}$. Además, para garantizar la descorrelación estructural entre los árboles y evitar que aprendan patrones idénticos, en cada división de nodo (\textit{split}) el algoritmo evalúa únicamente un subconjunto aleatorio de características extraídas desde la matriz global de datos. Finalmente, la predicción del ensamble para un problema de clasificación binaria ($Y \in {0,1}$) se obtiene formalmente mediante votación por mayoría, descrita en la Ecuación \ref{eq:rf_vote}:
\begin{equation}
    \hat{Y}(X) = \arg\max_{c \in \{0,1\}} \sum_{b=1}^{B} \mathbb{I}(C_b(X) = c)
    \label{eq:rf_vote}
\end{equation}

Donde $C_b(X)$ representa la predicción de clase emitida por el árbol número $b$ e $\mathbb{I}(\cdot)$ corresponde a la función indicadora, la cual toma el valor de 1 cuando la condición interna se cumple y 0 en caso contrario. La naturaleza no paramétrica de este algoritmo le permite modelar interacciones altamente no lineales sin requerir transformaciones previas de los datos, además de presentar una elevada resistencia al sobreajuste (\textit{overfitting}) debido al promedio estocástico de sus componentes (\cite{hastie2009elements}).

\textbf{XGBoost (\textit{eXtreme Gradient Boosting})}, a diferencia del enfoque basado en \textit{bagging}, utiliza la técnica de \textit{boosting}, la cual consiste en construir árboles de decisión de manera secuencial y aditiva. Cada nuevo árbol, denominado aprendiz débil, se optimiza específicamente para corregir los errores residuales (gradientes) calculados en las etapas anteriores; es decir, los errores cometidos por el conjunto previo de árboles. Matemáticamente, XGBoost minimiza directamente una función objetivo regularizada $\mathcal{L}$ en cada iteración $k$, equilibrando el ajuste empírico con la complejidad estructural de cada árbol $f_k$, tal como se expresa en la Ecuación \ref{eq:xgboost}:

\begin{equation}
    \mathcal{L} = \sum_{i=1}^{n} l(y_i, \hat{y}_i) + \sum_{k=1}^{K} \Omega(f_k)
    \label{eq:xgboost}
\end{equation}

Donde $l(\cdot)$ representa la función de pérdida diferenciable (por ejemplo, entropía cruzada para clasificación binaria), encargada de medir la discrepancia entre la etiqueta real $y_i$ y la estimación $\hat{y}_i$. Por otro lado, el término diferenciador clave del algoritmo corresponde al penalizador de complejidad topológica $\Omega(f_k)$, el cual restringe explícitamente el tamaño de la estructura para mitigar el sobreajuste. Este término de regularización se define mediante la Ecuación \ref{eq:xgboost_reg}:

En donde $l(\cdot)$ representa la función de pérdida diferenciable (como la entropía cruzada para variables binarias) que mide la discrepancia entre la etiqueta real $y_i$ y la estimación $\hat{y}_i$. Por otro lado, el término diferenciador clave del algoritmo es el penalizador de complejidad topológica $\Omega(f_k)$, el cual restringe explícitamente el tamaño de la estructura para mitigar el sobreajuste. Este término de regularización se define específicamente en la Ecuación \ref{eq:xgboost_reg} (\cite{chen2016xgboost}):

\begin{equation}
    \Omega(f_k) = \gamma T + \frac{1}{2}\lambda \sum_{j=1}^{T} w_j^2
    \label{eq:xgboost_reg}
\end{equation}

En esta expresión, $T$ corresponde al número total de hojas presentes en el árbol $f_k$, mientras que $w_j$ representa el peso asociado a cada hoja $j$. El hiperparámetro $\gamma$ actúa como una penalización sobre el número de divisiones del árbol, funcionando como un umbral mínimo de ganancia requerido para permitir una nueva partición. En paralelo, el parámetro $\lambda$ aplica regularización L2 (\textit{Ridge}) sobre los pesos de las hojas, suavizando las actualizaciones numéricas individuales. Esta combinación matemática otorga a XGBoost una elevada estabilidad computacional en estructuras tabulares dispersas y una capacidad nativa para manejar desbalances de clases mediante ajuste de gradientes (\cite{chen2016xgboost}).

\subsection{Inteligencia artificial explicable (XAI) y valores de Shapley (SHAP)}

Si bien los modelos de aprendizaje automático avanzados como \textit{Random Forest} y \textit{XGBoost} ofrecen altos rendimientos predictivos en la identificación de patrones no lineales, operan arquitectónicamente como modelos de caja negra (\textit{black-box}), lo que significa que permiten clasificar o predecir un desenlace (como el fallecimiento de pacientes), pero no explican de forma transparente cómo cada variable de entrada contribuye a esa predicción. Esta falta de interpretabilidad algorítmica dificulta la toma de decisiones clínicas y limita su aporte a la gestión sanitaria pública, donde la transparencia constituye un requisito ético y regulatorio. Para superar esta barrera, la investigación integra la metodología SHAP (\textit{SHapley Additive exPlanations}) \cite{lundberg2017shap}.

Este método fue desarrollado por Lundberg y Lee, y traduce los principios de la teoría de juegos cooperativos (introducida por Lloyd Shapley en 1953) al dominio del aprendizaje automático. En esta analogía matemática, la predicción del modelo equivale al pago total del juego, mientras que las características del paciente (como la edad, los procedimientos o los diagnósticos) actúan como jugadores. En este contexto, SHAP calcula la contribución marginal exacta $\phi_i$ de cada variable evaluando su impacto sobre todas las permutaciones posibles de características, según la Ecuación \ref{eq:shap}:
\begin{equation}
    \phi_i = \sum_{S \subseteq F \setminus \{i\}} \frac{|S|! (|F| - |S| - 1)!}{|F|!} [f_x(S \cup \{i\}) - f_x(S)]
    \label{eq:shap}
\end{equation}

En donde $F$ corresponde al conjunto de todas las características, $S$ representa un subconjunto de características que excluye a la variable $i$ (cuya contribución se desea calcular), y $f_x(S)$ corresponde a la predicción del modelo condicionada únicamente a las características presentes en el subconjunto $S$ \cite{lundberg2017shap}.


Si se aplica SHAP directamente sobre un problema de clasificación binaria (por ejemplo, fallecido y sobreviviente), el valor SHAP parte de un valor base correspondiente a la prevalencia promedio o probabilidad base estimada por el modelo. A partir de dicho punto, la metodología entrega una atribución direccional transparente:

\begin{itemize}
    \item \textbf{Valores SHAP positivos:} indican que la presencia (o una magnitud elevada) de una variable empuja la predicción del modelo a favor de la Clase 1, atribuyendo estadísticamente el riesgo hacia la presencia de esa clase (por ejemplo, mortalidad).
    \item \textbf{Valores SHAP negativos:} indican que la variable ejerce una contribución inversa, disminuyendo, por ejemplo, la probabilidad de fallecimiento y asociando el perfil clínico del paciente hacia la Clase 0 (pacientes sobrevivientes).
\end{itemize}

\section{Estado del arte}

\subsection{Estrategia de búsqueda y selección de estudios}

Para la revisión sistemática de la literatura se establecieron dos estrategias diferentes, una búsqueda global en \textbf{Scopus} para identificar el estado del arte internacional y una búsqueda local en Chile en \textbf{SciELO}.


En primer lugar, se ejecutó una consulta en Scopus para responder las siguientes preguntas de investigación sobre variables, tecnologías y metodologías utilizadas en la caracterización de pacientes oncológicos:

\begin{itemize}
    \item \textbf{Variables y factores asociados:} ¿Qué atributos o variables hospitalarias, administrativas y clínicas (registradas en GRD u otros sistemas de datos hospitalarios) se han utilizado en la literatura para caracterizar pacientes con cáncer, tanto en Chile como a nivel internacional?
    
    \item \textbf{Tecnologías:} ¿Qué enfoques computacionales (estadística, minería de datos y aprendizaje automático) se han aplicado para caracterizar perfiles de pacientes oncológicos utilizando registros clínicos y administrativos, como datos hospitalarios (por ejemplo, GRD), registros electrónicos o seguros de salud?
    
    \item \textbf{Metodologías exitosas o existentes:} ¿Qué metodologías han demostrado ser más efectivas para identificar patrones y atributos relevantes en la caracterización del cáncer a partir de registros clínicos u hospitalarios?
\end{itemize}


Dicha consulta fue la siguiente:

\begin{lstlisting}[style=query]
("cancer" OR "oncology" OR "neoplasms")
AND
("Diagnosis Related Groups" OR DRG OR "hospital discharge data" 
OR "administrative hospital data" OR "hospital data" 
OR "electronic health records" OR EHR OR "claims data")
AND
(attributes OR features OR variables OR "risk factors" 
OR "clinical profile" OR "patient profiling" 
OR "patient characterization" OR "clinical characterization")
AND
("machine learning" OR "data mining" OR "artificial intelligence" 
OR "statistical analysis" OR "predictive modeling" 
OR "risk prediction" OR "pattern recognition" OR "feature selection")
\end{lstlisting}


De los 182 resultados iniciales, se aplicaron los siguientes filtros, reduciendo la búsqueda a 104 resultados:
\begin{itemize}
    \item \textbf{Año:} 2020--2025
    \item \textbf{Área temática:} Medicina, Informática, Multidisciplinaria y Profesionales de la Salud
    \item \textbf{Tipo de documento:} Artículo y revista
    \item \textbf{Idioma:} Inglés y español
\end{itemize}

Posteriormente, se aplicó un proceso de selección basado en tres filtros secuenciales para identificar los artículos más relevantes para el objetivo de esta propuesta, a través de la lectura de sus títulos y resúmenes:

\begin{enumerate}
    \item \textbf{Filtro de relevancia central:} Se seleccionaron principalmente artículos que utilizan aprendizaje automático supervisado en oncología con datos análogos a GRD (como datos EHR y administrativos), y se descartaron artículos cuyo enfoque principal no correspondía al aprendizaje automático, como NLP puro o análisis de imágenes y radiomía.

    \item \textbf{Filtro de alineación de objetivos:} Se priorizaron estudios que predicen resultados similares a los de esta propuesta, como mortalidad, supervivencia, severidad, complicaciones y consumo de recursos o costos.
    
    \item \textbf{Filtro de relevancia metodológica:} Se seleccionaron artículos que utilizan técnicas de interpretabilidad como SHAP para crear perfiles de riesgo.
\end{enumerate}

Finalmente, este proceso resultó en la selección de 12 artículos internacionales, los cuales se resumen más adelante en la Tabla \ref{tab:resumen_scopus}.

En segundo lugar, se ejecutó una búsqueda en SciELO orientada a identificar estudios sobre la aplicación de aprendizaje automático sobre datos GRD en el área de oncología; sin embargo, esta consulta arrojó 0 resultados.

Entonces, con el fin de adoptar un enfoque descriptivo y responder a la necesidad de construir una línea base contextual sólida, se ejecutó una segunda búsqueda ampliada orientada a identificar estudios que, mediante estadística descriptiva o análisis epidemiológico, caracterizaran el cáncer en la población chilena utilizando datos locales. Para ello, se ejecutó la siguiente consulta:

\begin{lstlisting}[style=query]
(TITLE:("cancer" OR "oncologia" OR "tumor" OR "neoplasia"))
AND (clínico OR demográfico OR epidemiológico OR "egreso hospitalario" OR "carga de enfermedad")
AND (Chile OR "población chilena" OR "hospital chileno")
\end{lstlisting}

De 1583 resultados iniciales, se aplicaron los siguientes filtros, reduciendo la búsqueda a 44 resultados:

\begin{itemize}
    \item \textbf{Año:} 2018--2025
    \item \textbf{País:} Chile, para captar estudios locales.
\end{itemize}

Tras aplicar los filtros anteriores, se leyeron los resúmenes de los artículos y se aplicaron los siguientes criterios de selección final para asegurar su relevancia con el problema de los GRD:

\begin{itemize}
    \item \textbf{Enfoque poblacional o de cohorte:} Se seleccionaron estudios con muestras significativas (N $>$ 100) o estadísticas nacionales (DEIS), descartando reportes de casos únicos que no permiten generalizar perfiles.
    \item \textbf{Contexto local y territorial:} Se incluyeron investigaciones de hospitales públicos chilenos (Temuco y Chillán) para validar variables territoriales, como la ruralidad y el tipo de establecimiento.
    \item \textbf{Variables presentes en GRD:} Se seleccionaron estudios que analizaron variables presentes en GRD, como edad, sexo y mortalidad.
\end{itemize}

Finalmente, en la búsqueda de SciELO se seleccionaron 4 artículos fundamentales, los cuales constituyen la evidencia base sobre la carga demográfica y hospitalaria del cáncer en Chile, indicados más adelante en la Tabla \ref{tab:resumen_scielo}.


\subsection{Estudios sobre caracterización de cohortes oncológicas mediante aprendizaje automático e inteligencia artificial explicable}

En la literatura se identifican dos principales corrientes metodológicas en el uso de datos sanitarios en oncología. La primera corresponde a la predicción clínica mediante aprendizaje automático e inteligencia artificial, donde se ha demostrado que los registros clínicos y administrativos, análogos a los GRD (como los Registros de Salud Electrónicos, EHR), contienen patrones latentes capaces de predecir desenlaces como mortalidad, supervivencia, duración de la hospitalización y readmisión, utilizando algoritmos como Random Forest, XGBoost y regresión LASSO (\cite{article2}; \cite{article3}; \cite{article10}).


Para enfrentar el desbalance de clases, se aplicaron técnicas de remuestreo, como submuestreo o SMOTE, que consiste en crear muestras sintéticas de la clase minoritaria interpolando entre ejemplos existentes (\cite{article4}; \cite{article6}; \cite{article16}). Además, se emplearon validaciones temporales o prospectivas, separando el conjunto de entrenamiento y prueba por fecha para evitar fugas de datos (uso de información en el entrenamiento que no estaría disponible en el momento de la predicción). La calibración de los modelos se evaluó mediante la puntuación de Brier, que cuantifica la diferencia cuadrática media entre las probabilidades predichas y los desenlaces observados (\cite{article1}; \cite{article8}; \cite{article9}).

Asimismo, se incorporaron técnicas de interpretabilidad como SHAP, estableciendo un estándar de explicabilidad dual: global, para caracterizar factores de riesgo a nivel poblacional, y local, para justificar decisiones individuales, cuantificando la influencia de cada variable en las predicciones (\cite{article1}; \cite{article4}; \cite{article7}; \cite{article11}). Finalmente, estos modelos alcanzaron desempeños moderados y altos, con valores de AUC-ROC entre 0.71 y 0.94, evidenciando una capacidad predictiva superior al azar.

En específico, para esta corriente se tienen los siguientes estudios:
\begin{itemize}
\item \textbf{Desarrollo y validación de un modelo de aprendizaje automático explicable para predecir el pronóstico en pacientes con sepsis y antecedentes de cáncer ingresados en la unidad de cuidados intensivos:} El objetivo de este estudio fue construir y validar un modelo de aprendizaje automático explicable para predecir la mortalidad intrahospitalaria en pacientes con sepsis y antecedentes de cáncer ingresados en la unidad de cuidados intensivos (UCI). Para ello, se realizó un análisis retrospectivo de 3364 pacientes presentes en los registros de salud electrónicos (EHR) de la base de datos MIMIC-IV, que contiene información de pacientes ingresados en la UCI del hospital Beth Israel Deaconess Medical Center. Como metodología, la variable a predecir fue la mortalidad hospitalaria (como variable binaria); se seleccionaron 11 variables pronósticas clave mediante regresión LASSO y eliminación recursiva de características (RFE); se entrenaron ocho algoritmos de aprendizaje automático, como Random Forest (RF), XGBoost, árbol de decisión (DT) y SVM; y se aplicaron técnicas de interpretabilidad mediante SHAP y gráficos de dependencia parcial (PDPs). Finalmente, Random Forest fue el modelo con mayor precisión predictiva, con un AUC de 0.78 y el menor Brier Score. Además, el análisis de explicabilidad reveló que la característica más influyente para la mortalidad fue la puntuación de gravedad APSIII (que mide la disfunción fisiológica aguda III), junto con la puntuación SOFA, la saturación de oxígeno (SpO$_2$), y la frecuencia cardíaca y respiratoria (\cite{article1}).
\item \textbf{Desarrollo de un modelo de IA para predecir el estado de hipoxia y el pronóstico en el cáncer de pulmón de células no pequeñas utilizando datos multimodales:} El objetivo de este estudio fue crear un modelo de inteligencia artificial (IA) para pacientes con cáncer de pulmón de células no pequeñas (NSCLC) que integrara datos multimodales para predecir el estado de hipoxia y la supervivencia general. Como metodología, se utilizaron datos de 452 pacientes provenientes de múltiples cohortes clínicas y bases públicas. En específico, se emplearon datos multimodales que incluyeron registros de salud electrónicos (EHR), características radiómicas de imágenes de tomografía computarizada (TC) y conjuntos de datos genómicos como The Cancer Genome Atlas (TCGA) y Gene Expression Omnibus (GEO). Para la construcción del modelo, los pacientes se clasificaron en grupos de hipoxia alta o baja en función de su expresión génica mediante el método ssGSEA. Posteriormente, se entrenaron modelos específicos para cada tipo de dato: LASSO para las características radiómicas de TC, ssuBERT para extraer información fenotípica estructurada desde textos libres de EHR y una red neuronal profunda (DNN) para integrar las características provenientes de las distintas fuentes de datos. Finalmente, el modelo multimodal obtuvo un AUC promedio de 0.8449, demostrando que la integración transversal de diferentes arquitecturas de datos mejoró sustancialmente la predicción pronóstica frente a enfoques unimodales. Además, se concluyó que niveles elevados de hipoxia se correlacionaban con peores resultados de supervivencia global, y el modelo unimodal con mejor desempeño fue ssuBERT, con un AUC de 0.7662 (\cite{article2}).
\item \textbf{Hacia la predicción de la rehospitalización a los 30 días en pacientes oncológicos: identificación de factores de riesgo oportunos y accionables:} El objetivo de este estudio fue desarrollar modelos de predicción de readmisión hospitalaria no planificada a 30 días para pacientes oncológicos (variable binaria), con el fin de identificar intervenciones oportunas y accionables. Para ello, se utilizaron datos de registros de salud electrónicos (EHR) de una cohorte retrospectiva de 2460 pacientes oncológicos del centro de tratamiento Ann B. Barshinger Cancer Institute (ABBCI). Como metodología, se estableció un pipeline de trabajo que incluyó el preprocesamiento, la selección de características y el entrenamiento de 11 algoritmos de aprendizaje automático, como Random Forest, XGBoost y LGB; además, se utilizó Optuna (optimización bayesiana) para optimizar los hiperparámetros de los modelos, junto con métricas de evaluación como ROC y AUPRC, validación cruzada anidada de 10 pliegues y SHAP para la identificación de los factores de riesgo más importantes para las predicciones. Cabe destacar que los múltiples clasificadores de ML se entrenaron comparando el rendimiento al utilizar todas las variables de la visita clínica frente a utilizar solo aquellas disponibles durante las primeras 48 horas de admisión. Finalmente, el modelo Light Gradient Boosting (LGB) obtuvo el mejor rendimiento global (con un AUROC de 0.711) utilizando todas las características, mientras que Random Forest presentó el mejor desempeño utilizando datos limitados a las primeras 48 horas (con un AUROC de 0.684). Dentro de los predictores más influyentes de la readmisión hospitalaria se encontraron factores de riesgo dinámicos, como albúmina baja y frecuencia cardiaca elevada al alta, además del cambio de peso anual, síntomas de depresión y el nivel de ingresos según código postal. Asimismo, se demostró que el aprendizaje automático puede igualar o superar modelos propietarios existentes y ofrecer perspectivas clínicas y socioeconómicas accionables de forma temprana (\cite{article3}).
\item \textbf{Predicción dinámica de metástasis en ganglios linfáticos laterales mediante aprendizaje automático en pacientes con cáncer papilar de tiroides:} El objetivo de este estudio fue desarrollar un modelo de aprendizaje automático basado en la web para predecir dinámicamente la metástasis en ganglios linfáticos laterales (LLNM) en pacientes con cáncer papilar de tiroides (PTC) antes de la cirugía. Para ello, se analizaron retrospectivamente datos de ultrasonido y EHR clínicos de un centro médico en China (First Medical Centre of the Chinese PLA General Hospital), con una cohorte de 1815 pacientes. Como metodología, se aplicó regresión LASSO para seleccionar las 30 características más relevantes para construir los modelos; se entrenaron 6 algoritmos de aprendizaje automático, incluyendo XGBoost, Random Forest, árbol de decisión, SVM, red neuronal y KNN, y se compararon con un nomograma de regresión logística tradicional. Cada modelo fue evaluado en función de su precisión, sensibilidad, AUC y DCA; además, se utilizó SHAP para identificar las características más influyentes en las predicciones del modelo. Finalmente, el modelo Random Forest demostró ser el más robusto, con un AUC de 0.8 y una sensibilidad de 0.69. SHAP identificó que las características más influyentes para predecir LLNM fueron el tamaño del tumor, la edad, la microcalcificación, el tamaño de los ganglios linfáticos y la ubicación del tumor. Finalmente, la investigación entregó una herramienta accesible para facilitar la estratificación de riesgos a nivel individual (\cite{article4}).
\item \textbf{Comparación de la regresión logística con algoritmos de aprendizaje automático para la predicción del riesgo de cáncer gástrico en flujos de datos clínicos reales:} El objetivo de este estudio fue comparar la eficacia de la regresión logística (LR) frente a algoritmos de aprendizaje automático (ML) para predecir el riesgo de cáncer gástrico no cardial (NCGC) utilizando datos clínicos reales. Para ello, se analizaron retrospectivamente datos de registros de salud electrónicos (EHR) de dos sistemas de salud académicos en Estados Unidos: Stanford y la Universidad de Washington (UW), con una cohorte total de 59,000 pacientes. Como metodología, se comparó el modelo clásico de regresión logística (basado en factores de riesgo preestablecidos como histología, infección previa por \textit{Helicobacter pylori}, raza, etnia, tabaquismo y anemia) con cuatro algoritmos de aprendizaje automático: Random Forest, SVM, LASSO penalizado y KNN (que seleccionaron variables de manera agnóstica), utilizando métricas de evaluación como exactitud, sensibilidad, especificidad y curvas AUC/ROC. Finalmente, el rendimiento del modelo de regresión logística no fue inferior al de los modelos de ML en la validación externa, logrando una precisión del 72.4\%. Además, en validaciones tanto internas como externas, la regresión logística mostró una precisión, sensibilidad y especificidad comparables a los modelos de aprendizaje automático optimizados. Por ende, el estudio concluyó que los modelos estadísticos clásicos construidos sobre variables predictoras robustas y seleccionadas manualmente no son inferiores a los modelos de aprendizaje automático impulsados por datos masivos para la predicción del riesgo de cáncer gástrico. Por otra parte, dentro de los modelos de ML, el modelo KNN mostró una mayor precisión y especificidad (\cite{article5}).
\item \textbf{Factores de riesgo asociados con eventos relacionados con el esqueleto tras la interrupción del tratamiento con denosumab en pacientes con metástasis óseas de tumores sólidos: un enfoque de aprendizaje automático en el mundo real:} El objetivo de este estudio fue identificar los factores de riesgo asociados a la aparición de eventos relacionados con el esqueleto (SREs) tras la interrupción del tratamiento con denosumab (un anticuerpo monoclonal) en pacientes con metástasis óseas derivadas de tumores sólidos. Para ello, se analizaron retrospectivamente datos de registros de salud electrónicos (EHR) de Optum PanTher, considerando pacientes diagnosticados entre enero de 2007 y septiembre de 2019. Como metodología, se utilizó el algoritmo XGBoost para las predicciones; se gestionó el desbalance de clases mediante técnicas de sobremuestreo, submuestreo y penalización. Además, se ajustaron los hiperparámetros mediante optimización bayesiana con validación cruzada de 5 pliegues, y el rendimiento del modelo se evaluó mediante las métricas de precisión, sensibilidad, AUROC y especificidad. Para la interpretabilidad, se aplicó SHAP para evaluar sin sesgos la importancia de más de 60 variables predictoras. Como resultado, se obtuvo que el modelo XGBoost alcanzó un AUROC de 77\% en la predicción de SREs, y dentro de los factores de riesgo más influyentes para la detección de SREs se encontraron los SREs previos, una menor duración del tratamiento con denosumab, la edad joven y la presencia de cáncer de próstata. Finalmente, este estudio demostró que el aprendizaje automático puede guiar las decisiones sobre la continuidad del tratamiento con denosumab para mejorar los desenlaces del paciente (\cite{article6}).
\item \textbf{Aprendizaje automático para la clasificación del estado postoperatorio del paciente utilizando datos médicos estandarizados:} Este estudio tuvo como objetivo identificar factores de riesgo e interacciones temporales que contribuyen a una hospitalización postoperatoria prolongada en pacientes con cáncer de pulmón. Para ello, el estudio empleó Datos del Mundo Real (RWD) estandarizados, provenientes de los sistemas Diagnosis Procedure Combination (DPC) y vías clínicas (Clinical Pathway o CP). Como metodología, se utilizaron métodos de agrupamiento o \textit{clustering} (algoritmo CLUTO) para agrupar pacientes según patrones de ítems DPC, en conjunto con la construcción de modelos de aprendizaje automático. Además, se aplicó la técnica SHAP para cuantificar la contribución de cada variable a la predicción de hospitalización prolongada, junto con un gráfico dirigido para explorar las relaciones temporales entre las variables que conducen a hospitalizaciones largas en vez de cortas. Como resultados, el modelo identificó que el factor de riesgo más influyente en la hospitalización prolongada fue la varianza de estreñimiento (en POD1-4), y el \textit{clustering} agrupó variables como tabaquismo, obesidad, estado respiratorio y estreñimiento, por un lado, y el manejo de drenajes, heridas y goteo, por otro. Finalmente, este estudio demostró que la minería de datos sobre registros médicos estandarizados permite revelar vías causales que apoyan la toma de decisiones clínicas tempranas (\cite{article7}).
\item \textbf{Aprendizaje automático aplicado a los registros electrónicos de salud: identificación de pacientes de quimioterapia con alto riesgo de visitas evitables a urgencias e ingresos hospitalarios:} Este estudio tuvo como objetivo predecir el riesgo de uso prevenible de cuidados agudos (ACU), como visitas a urgencias o ingresos hospitalarios no planificados, en pacientes oncológicos que están iniciando quimioterapia. Para ello, se analizaron datos estructurados de registros de salud electrónicos (EHR) generados antes del inicio de la quimioterapia en pacientes adultos tratados en sitios comunitarios afiliados entre enero de 2013 y julio de 2019. Como metodología, se entrenaron, con información generada antes de iniciar el tratamiento, nueve modelos de aprendizaje automático, incluyendo LASSO, Random Forest, Ridge, Elastic Net, redes neuronales perceptrón multicapa (MLP), LightGBM, SVM, KNN y un modelo de votación ensamblado. Como resultado, el modelo LASSO obtuvo el mejor rendimiento predictivo a 180 días, con un AUROC de 0.783. Además, dentro de las variables clave seleccionadas por el modelo LASSO se encontraban hospitalizaciones previas, diagnóstico de depresión, saturación de oxígeno y pulso, entre otras. Finalmente, este estudio validó que los modelos de aprendizaje automático basados en EHR integrales pueden identificar eficazmente a los pacientes en riesgo para focalizar intervenciones preventivas (\cite{article8}).
\item \textbf{Rendimiento de un algoritmo de aprendizaje automático que utiliza datos de registros médicos electrónicos para identificar y estimar la supervivencia en una cohorte longitudinal de pacientes con cáncer de pulmón:} El objetivo de este estudio consistió en investigar la viabilidad de ensamblar una cohorte clínica longitudinal de cáncer de pulmón a partir de datos de registros médicos electrónicos (EHR) mediante aprendizaje automático y utilizarla para desarrollar un modelo pronóstico de supervivencia. Para ello, se utilizaron registros de más de 76,000 pacientes del sistema de salud Mass General Brigham (MGB) desde julio de 1988 hasta octubre de 2018. Como metodología, se utilizó el algoritmo PheCAP para fenotipado, empleando dentro de este los métodos SAFE (para la selección de características) y LASSO (para la clasificación principal a partir de las características seleccionadas por SAFE). Además, se aplicaron herramientas de Procesamiento de Lenguaje Natural (NLP) para extraer datos clínicos de textos libres escritos por médicos (datos no estructurados). En específico, se utilizaron los algoritmos NICE (desarrollado por los propios investigadores) y EXTEND (para extraer datos numéricos específicos, como el índice de masa corporal, IMC). Asimismo, para predecir la supervivencia de los pacientes, se utilizaron algoritmos de regresión de Cox penalizada y penalización MCP para seleccionar solo las variables más importantes en la predicción de la supervivencia global a 1, 2, 3, 4 y 5 años. Como resultados, el algoritmo identificó a 42,069 pacientes con cáncer de pulmón, logrando un valor predictivo positivo (VPP) del 94.4\% y un AUROC de 0.927. Además, el modelo pronóstico de supervivencia basado en Cox obtuvo un AUROC de 0.828 para la predicción a 1 año y de 0.812 para la predicción a 5 años. Finalmente, este estudio demostró que los EHR proporcionan una alta fidelidad para extraer información compleja y predecir resultados a largo plazo (\cite{article9}).
\item \textbf{Modelo de predicción de datos clínicos para identificar pacientes con cáncer de páncreas en etapa temprana:} El objetivo de este estudio fue evaluar un modelo de aprendizaje automático para identificar pacientes con cáncer de páncreas en etapas tempranas (ESPC), basándose exclusivamente en datos clínicos rutinarios alojados en registros médicos electrónicos (EHR). Para ello, se extrajeron datos de la base de datos Optum (de Estados Unidos), la cual contiene registros diagnósticos de pacientes entre 2009 y 2017, donde 3322 presentaban cáncer de páncreas en etapa temprana (ESPC), definidos por haber recibido cirugía pancreática mayor, y 25,908 presentaban cáncer de páncreas en etapa tardía (LSPC). Como metodología, se construyó y optimizó un modelo predictivo utilizando XGBoost, y se ajustaron los hiperparámetros mediante \textit{Grid Search} y validación cruzada de 5 pliegues. Como resultado, el modelo XGBoost logró distinguir los casos con un AUC de 0.84. Además, el algoritmo demostró la capacidad de identificar el 58\% de los casos de cáncer de páncreas diagnosticados en etapa tardía (LSPC) hasta 24 meses antes del diagnóstico clínico formal, con un 60\% de sensibilidad, lo que resalta el potencial del aprendizaje automático para el diagnóstico temprano y la reducción de la mortalidad en neoplasias agresivas como el cáncer de páncreas (\cite{article10}).
\item \textbf{Un modelo multicéntrico de bosque aleatorio para la predicción eficaz del pronóstico en una red de investigación clínica colaborativa:} El objetivo de este estudio fue proponer un modelo de aprendizaje automático basado en Random Forest (RF) multicéntrico para la predicción de pronósticos (supervivencia a 5 años) en pacientes con cáncer colorrectal (CRC), permitiendo la minería de datos clínicos federados a partir de bases de datos particionadas horizontalmente, garantizando la privacidad de los pacientes sin necesidad de compartir datos crudos y sensibles. Para ello, se utilizaron registros provenientes de la base de datos SEER (de Estados Unidos) y datos del hospital SAHZU (de China). Como metodología, el estudio desarrolló un enfoque de mejora de datos basado en Redes Generativas Antagónicas con privacidad diferencial (DPGAN) para generar datos sintéticos locales. Posteriormente, se evaluaron los modelos bajo diferentes niveles de heterogeneidad institucional y se comparó el modelo multicéntrico de Random Forest con uno entrenado de forma centralizada y con otros clasificadores clásicos, como regresión logística, SVM y redes neuronales. Adicionalmente, se comparó con un modelo de estadificación TNM como línea base clínica; se implementó un método de ranking de importancia de variables basado en la impureza de Gini y se utilizaron métricas de evaluación como AUC (ROC), exactitud, sensibilidad, especificidad y AUPRC. Como resultados, el modelo multicéntrico de Random Forest obtuvo un mejor rendimiento tanto en la validación cruzada interna (AUC de 0.875 con los datos SEER) como en la validación externa (AUC de 0.852 con los datos de SAHZU), mostrando un desempeño superior al de la regresión logística (AUC de 0.835) y al modelo de estadificación TNM (AUC de 0.771). Finalmente, el procedimiento de selección de características demostró ser efectivo para simplificar el modelo sin perder exactitud, concluyendo que este enfoque ofrece una técnica práctica y segura para construir modelos pronósticos robustos en redes de investigación oncológica colaborativas (\cite{article11}).
\item \textbf{Un marco de aprendizaje automático explicable para la predicción de la duración de la estancia hospitalaria en pacientes con cáncer de pulmón:} El objetivo de este estudio fue desarrollar un marco de trabajo de aprendizaje automático explicable para predecir el tiempo de estancia hospitalaria (LOS) de pacientes con cáncer de pulmón ingresados en la unidad de cuidados intensivos (UCI), abordando el desafío del desbalance de clases en los registros de salud electrónicos (EHR). Para ello, se analizó la base de datos pública Medical Information Mart for Intensive Care (MIMIC-III), proveniente de pacientes adultos que estuvieron en la UCI del hospital Beth Israel Deaconess Medical Center (BIDMC), en Massachusetts, Estados Unidos, entre 2001 y 2012. Como metodología, se estructuró el problema predictivo como una clasificación binaria (estancia corta versus estancia larga) y se entrenaron modelos como Random Forest (RF), XGBoost y regresión logística. Para optimizar el entrenamiento, se aplicaron métodos de selección de variables (significancia clínica y eliminación recursiva de características, RFE) y se evaluaron seis técnicas de balanceo de clases, incluyendo sobremuestreo (SMOTE, ADASYN), submuestreo (ENN, TomekLinks) y métodos combinados (SMOTE-ENN, SMOTETomek). Para la evaluación de los modelos se utilizó validación cruzada de 10 iteraciones (\textit{10-fold cross-validation}) y métricas de sensibilidad, especificidad, AUC e \textit{Index Balanced Accuracy} (IBA). Por último, se aplicó SHAP para la interpretabilidad clínica de las predicciones. Como resultado, el clasificador Random Forest demostró un rendimiento superior y resistente a variaciones en la selección de características. Las técnicas de sobremuestreo obtuvieron los mejores resultados predictivos, donde ADASYN y SMOTE alcanzaron un Área Bajo la Curva (AUC) de 100\% y 98\%, respectivamente. Además, el análisis SHAP identificó que las características más influyentes para predecir la estancia hospitalaria fueron la temperatura, la admisión por emergencia, la glucosa, las plaquetas y la frecuencia respiratoria. Finalmente, se concluyó que la integración del algoritmo SMOTE con Random Forest y SHAP constituye una herramienta valiosa para apoyar a gestores de camas en la optimización de los recursos de la UCI (\cite{article16}).
\end{itemize}

Cabe mencionar que, en la Tabla \ref{tab:resumen_scopus} se presentan los resumenes de cada uno de los artículos internacionales seleccionados sobre la caracterización de cohortes oncológicas mediante aprendizaje automático e inteligencia artificial explicable.

\begin{table}[p]
\centering
\scriptsize
\setlength{\tabcolsep}{2pt}
\renewcommand{\arraystretch}{1.1}
\caption{Resumen de estudios sobre caracterización de cohortes oncológicas mediante aprendizaje automático e inteligencia artificial explicable.}
\label{tab:resumen_scopus}
\begin{tabular}{|c|c|c|c|}
\hline
\parbox[t]{1.2cm}{\centering\textbf{Estudio}} & \parbox[t]{4.0cm}{\centering\textbf{Objetivo principal}} & \parbox[t]{5.0cm}{\centering\textbf{Datos y metodología}} & \parbox[t]{5.0cm}{\centering\textbf{Hallazgos principales}} \\ \hline
\parbox[t]{1.2cm}{\centering \cite{article1}} & \parbox[t]{4.0cm}{Predecir la mortalidad intrahospitalaria en pacientes con sepsis y antecedentes de cáncer en la UCI.} & \parbox[t]{5.0cm}{\textbf{Datos:} MIMIC-IV (EHR). \newline \textbf{Métodos:} LASSO, RFE, 8 modelos ML (incluyendo RF y XGBoost), SHAP, PDPs.} & \parbox[t]{5.0cm}{RF obtuvo la mayor precisión (AUC 0.78). SHAP identificó las puntuaciones APSIII, SOFA y SpO$_2$ como los predictores más influyentes.} \\ \hline
\parbox[t]{1.2cm}{\centering \cite{article2}} & \parbox[t]{4.0cm}{Predecir el estado de hipoxia y la supervivencia global en cáncer de pulmón (NSCLC).} & \parbox[t]{5.0cm}{\textbf{Datos:} Multimodales (EHR, TC radiómica, TCGA, GEO). \newline \textbf{Métodos:} ssGSEA, LASSO, ssuBERT, DNN (Red Neuronal Profunda)} & \parbox[t]{5.0cm}{El modelo multimodal superó a los unimodales (AUC 0.8449). Altos niveles de hipoxia se correlacionaron con una peor supervivencia.} \\ \hline
\parbox[t]{1.2cm}{\centering \cite{article3}} & \parbox[t]{4.0cm}{Identificar factores de riesgo para predecir la readmisión hospitalaria a 30 días en oncología.} & \parbox[t]{5.0cm}{\textbf{Datos:} EHR (ABBCI). \newline \textbf{Métodos:} 11 algoritmos ML, Optuna, validación cruzada anidada, SHAP.} & \parbox[t]{5.0cm}{LGB fue el mejor con datos completos (AUROC 0.711); RF destacó con datos de 48h (AUROC 0.684). SHAP: Albúmina baja, pulso y depresión como claves.} \\ \hline
\parbox[t]{1.2cm}{\centering \cite{article4}} & \parbox[t]{4.0cm}{Predecir dinámicamente la metástasis en ganglios linfáticos laterales (LLNM) en cáncer papilar de tiroides.} & \parbox[t]{5.0cm}{\textbf{Datos:} Ultrasonido y EHR clínicos (China). \newline \textbf{Métodos:} LASSO, 6 algoritmos ML, nomograma logístico, SHAP. } & \parbox[t]{5.0cm}{F fue el modelo más robusto (AUC 0.80). SHAP reveló que el tamaño del tumor, edad y microcalcificación son factores críticos.} \\ \hline
\parbox[t]{1.2cm}{\centering \cite{article5}} & \parbox[t]{4.0cm}{Comparar la regresión logística (LR) frente al ML para predecir el riesgo de cáncer gástrico.} & \parbox[t]{5.0cm}{\textbf{Datos:} EHR (Stanford, UW). \newline \textbf{Métodos:} LR tradicional vs. ML agnóstico (RF, SVM, LASSO, KNN).} & \parbox[t]{5.0cm}{LR no fue inferior al ML (precisión 72.4\%), demostrando que los modelos clásicos con variables bien seleccionadas son altamente competitivos.} \\ \hline
\parbox[t]{1.2cm}{\centering \cite{article6}} & \parbox[t]{4.0cm}{Identificar riesgos de eventos relacionados con el esqueleto (SREs) tras interrumpir denosumab.} & \parbox[t]{5.0cm}{\textbf{Datos:} EHR (Optum PanTher). \newline \textbf{Métodos:} XGBoost, optimización bayesiana, técnicas de remuestreo, SHAP.} & \parbox[t]{5.0cm}{XGBoost alcanzó un AUROC del 77\%. Según SHAP: SREs previos, menor duración de tratamiento y cáncer de próstata aumentan el riesgo.} \\ \hline
\parbox[t]{1.2cm}{\centering \cite{article7}} & \parbox[t]{4.0cm}{Identificar interacciones que prolongan la estancia postoperatoria en cáncer de pulmón.} & \parbox[t]{5.0cm}{\textbf{Datos:} RWD estandarizados (DPC y vías clínicas). \newline \textbf{Métodos:} Clustering (CLUTO), ML, SHAP, gráficos dirigidos.} & \parbox[t]{5.0cm}{La varianza de estreñimiento fue el mayor riesgo de estancia prolongada. El clustering diferenció perfiles respiratorios vs. manejo de heridas.} \\ \hline
\parbox[t]{1.2cm}{\centering \cite{article8}} & \parbox[t]{4.0cm}{Predecir el uso prevenible de cuidados agudos (ACU) al iniciar quimioterapia.} & \parbox[t]{5.0cm}{\textbf{Datos:} EHR comunitarios. \newline \textbf{Métodos:} 9 modelos ML entrenados con datos previos al tratamiento.} & \parbox[t]{5.0cm}{LASSO obtuvo el mejor rendimiento (AUROC 0.783). Hospitalizaciones previas y diagnóstico de depresión fueron predictores clave.} \\ \hline
\parbox[t]{1.2cm}{\centering \cite{article9}} & \parbox[t]{4.0cm}{Ensamblar una cohorte longitudinal de cáncer de pulmón desde EHR para estimar supervivencia.} & \parbox[t]{5.0cm}{\textbf{Datos:} EHR (Mass General Brigham). \newline \textbf{Métodos:} Algoritmos PheCAP, NLP (NICE, EXTEND), regresión de Cox penalizada.} & \parbox[t]{5.0cm}{Identificación de cohorte con VPP del 94.4\%. El modelo de Cox estimó la supervivencia a 1 año con un excelente AUROC de 0.828.} \\ \hline
\parbox[t]{1.2cm}{\centering \cite{article10}} & \parbox[t]{4.0cm}{Identificar pacientes con cáncer de páncreas en etapas tempranas mediante datos rutinarios.} & \parbox[t]{5.0cm}{\textbf{Datos:} EHR (Optum). \newline \textbf{Métodos:} XGBoost, Grid Search, validación cruzada (5 pliegues).} & \parbox[t]{5.0cm}{XGBoost discriminó casos con un AUC de 0.84, identificando el 58\% de casos de etapa tardía hasta 24 meses antes del diagnóstico clínico. } \\ \hline
\parbox[t]{1.2cm}{\centering \cite{article11}} & \parbox[t]{4.0cm}{Predecir pronóstico de supervivencia a 5 años en cáncer colorrectal preservando la privacidad.} & \parbox[t]{5.0cm}{\textbf{Datos:} Bases descentralizadas (SEER, SAHZU). \newline \textbf{Métodos:} DPGAN (datos sintéticos), RF multicéntrico, Gini. } & \parbox[t]{5.0cm}{RF multicéntrico superó a LR y TNM (AUC 0.875) sin compartir datos crudos.} \\ \hline
\parbox[t]{1.2cm}{\centering \cite{article16}} & \parbox[t]{4.0cm}{Predecir el tiempo de estancia hospitalaria (LOS) en UCI para pacientes con cáncer de pulmón.} & \parbox[t]{5.0cm}{\textbf{Datos:} MIMIC-III. \newline \textbf{Métodos:} RF, XGBoost, LR, RFE, SMOTE/ADASYN, SHAP. } & \parbox[t]{5.0cm}{RF con sobremuestreo (ADASYN/SMOTE) alcanzó AUCs de 100\% y 98\%. SHAP identificó temperatura, glucosa y emergencia como factores críticos. } \\ \hline
\end{tabular}
\end{table}


\subsection{Caracterización de cohortes oncológicas mediante análisis epidemiológico y estadística descriptiva}
\label{sec:caracterizacion_cohortes_estadistica}

La segunda corriente corresponde a la caracterización epidemiológica y estadística descriptiva tradicional, que, mediante medidas de tendencia central, dispersión y distribuciones de frecuencia, permitió identificar grupos de riesgo en cohortes chilenas, como una mayor incidencia de cáncer de pulmón en hombres fumadores (\cite{scielo1}), el predominio de fallecimientos por cáncer oral y orofaríngeo en hombres mayores de 75 años (\cite{scielo2}), una mayor incidencia de cáncer de tiroides en mujeres (\cite{scielo3}) y un aumento del 109\% en muertes por cáncer entre 1986 y 2016, proyectándose como una potencial causa de muerte en Chile (\cite{scielo4}). Sin embargo, estos estudios se limitan a análisis descriptivos de centros médicos específicos y no modelan interacciones complejas (no lineales) entre variables, como la variabilidad en estancias hospitalarias y costos de atención (\cite{Wodchis2016}). Además, presentan sesgos ambientales, como la exposición a leña en la zona sur de Chile (Temuco) (\cite{scielo1}).

Por ende, la brecha identificada radica en la ausencia de metodologías que integren los registros oncológicos GRD con técnicas de aprendizaje automático supervisado explicable, capaces de transformar esta información en conocimiento accionable mediante perfiles que revelen los factores influyentes en las predicciones.

En específico, para esta corriente se tienen los siguientes estudios:

\begin{itemize}
\item \textbf{Perfil clínico, demográfico y anatomopatológico de pacientes con cáncer pulmonar en el Hospital de Temuco (2019-2023):} El objetivo de este estudio consistió en caracterizar el perfil clínico y demográfico de pacientes adultos con diagnóstico de cáncer pulmonar confirmado en el Hospital Regional de Temuco entre 2019 y 2023, y analizar su asociación con factores de riesgo locales (como tabaquismo y exposición a biomasa), síntomas al diagnóstico y presencia de mutaciones moleculares. Para ello, se analizó una cohorte retrospectiva extraída de fichas clínicas de 220 pacientes del Hospital Hernán Henríquez Aravena (Temuco), mayores de 15 años, con confirmación histológica (biopsia) de cáncer pulmonar, entre octubre de 2019 y junio de 2023. Como metodología, el diseño de este estudio fue observacional y descriptivo, donde se utilizó estadística descriptiva y se realizó un análisis exploratorio de datos, calculando medidas de tendencia central (promedio y mediana) y dispersión (desviación estándar) para variables numéricas (como la edad). Para las variables categóricas (como sexo o etapa), se utilizaron distribuciones de frecuencia absoluta y relativa (porcentajes). En específico, se evaluaron variables biosociodemográficas y clínicas (escala ECOG y etapa TNM), además de tipos de biopsia y estudios de patología molecular (EGFR, ALK y PDL-1). Como resultado, de los 220 pacientes confirmados con cáncer de pulmón, se obtuvo un perfil de 55.3\% de hombres con una edad promedio de 68 años (con una desviación estándar de 9.5 años), donde el 67.7\% tenía antecedentes de tabaquismo y un 29.5\% presentaba exposición a biomasa (leña). Además, se evidenció un diagnóstico alarmantemente tardío, con un 81.9\% de los casos detectados en etapas avanzadas (III-IV), siendo el adenocarcinoma el subtipo principal (63.9\%). Finalmente, se encontró una asociación significativa entre la expresión de PDL-1 y la exposición a biomasa, concluyendo que existe una necesidad urgente de implementar métodos de screening local (pruebas para detectar enfermedades o anomalías) para pacientes asintomáticos de alto riesgo (\cite{scielo1}).
\item \textbf{Mortalidad por cáncer oral y orofaríngeo en Chile entre los años 1955 al 2021:} El objetivo de este estudio fue determinar y analizar la variación histórica y la tendencia de las tasas de mortalidad por cáncer oral y orofaríngeo en la población chilena a lo largo de un período de 66 años (entre 1955 y 2021). Para ello, se utilizaron los registros oficiales de defunción del Instituto Nacional de Estadísticas (INE) y del Departamento de Estadísticas e Información en Salud (DEIS). Se calcularon las tasas brutas de mortalidad por 100000 habitantes y se estratificaron los hallazgos por género. En específico, este estudio fue cuantitativo, descriptivo y observacional, donde se ajustaron las tasas de mortalidad por edad (estandarizadas según la población mundial de la OMS) para hacer comparables los datos en el tiempo, y se utilizó regresión Joinpoint, la cual identifica matemáticamente puntos de quiebre donde la tendencia cambia significativamente de dirección, calculando el porcentaje de cambio anual (APC) para cada tramo: si es positivo, la mortalidad aumentó en ese período; si es negativo, disminuyó. Además, se usaron intervalos de confianza del 95\% para asegurar que las tendencias observadas no fueran producto del azar; es decir, si el estudio se repitiera 100 veces, en 95 de ellas el resultado estaría dentro de ese rango. Como resultado, de un total de 9345 muertes analizadas en el período estudiado, las tasas brutas de mortalidad oscilaron entre 0.92 y 1.53 por cada 100000 habitantes, evidenciando una clara tendencia al alza a nivel nacional, especialmente desde el año 2011. Además, la tasa de mortalidad se concentró fuertemente en el grupo de 75 años o más, y aunque los hombres presentaron una tasa de mortalidad mayor (1.64) en comparación con las mujeres (0.67), la brecha histórica de género ha disminuido. Finalmente, este estudio valida que las variables edad y sexo son predictores críticos para la mortalidad por cáncer oral y orofaríngeo, concluyendo que, a pesar de la leve reducción nacional en el consumo de tabaco, la mortalidad sigue aumentando, lo que evidencia la falta de políticas públicas efectivas para la detección y diagnóstico temprano de estas patologías específicas (\cite{scielo2}).
\item \textbf{Comportamiento del cáncer de tiroides en la Unidad de Cirugía Adulto del Hospital Clínico Herminda Martín de Chillán. Período 2017 al 2019:} El objetivo de este estudio fue describir el comportamiento clínico, demográfico e histopatológico del cáncer de tiroides en pacientes de la Región de Ñuble sometidos a resolución quirúrgica en la Unidad de Cirugía Adulto. Para ello, se analizaron los registros clínicos de 124 pacientes operados por patología tiroidea en el Hospital Clínico Herminda Martín de Chillán entre 2017 y 2019, de los cuales 58 fueron confirmados con patología maligna. Como metodología, se tabuló la correspondencia diagnóstica del sistema citológico de Bethesda, el tipo de resección (tiroidectomía) y parámetros de extensión tumoral. Para ello, se revisaron fichas clínicas, protocolos operatorios y biopsias de pacientes operados por patología tiroidea, calculando frecuencias relativas (porcentajes) para las variables cualitativas (como tipo de cirugía y sexo) y promedios para las variables cuantitativas (como edad y tiempo de espera). Como resultado, de los 124 pacientes operados, 58 presentaban cáncer (46.8\%), predominando las mujeres (86.2\%) y una edad promedio de 47.5 años. El diagnóstico más frecuente fue el carcinoma papilar (93.1\%) y el procedimiento más frecuente fue la tiroidectomía total (74.1\%), que implica la extirpación completa de la glándula tiroides. Además, clínicamente se reportó un 98\% de asociación simultánea entre la presencia de cáncer y cuadros de tiroiditis, y se midieron los tiempos de espera quirúrgica, encontrándose un promedio de 147 días para pacientes con cáncer. Finalmente, este estudio destaca que más de la mitad de los tumores femeninos medían más de 1 cm y que el 36\% ya presentaba invasión extratiroidea al momento de la cirugía, variables de extensión que los autores relacionan críticamente con los prolongados tiempos de espera en la red pública de salud. Cabe destacar que estos resultados son válidos para describir la realidad local del hospital, ya que se trató de un estudio censal de la unidad, considerando a todos los pacientes operados en el período, sin error de muestreo (\cite{scielo3}).
\item \textbf{Cáncer en Chile y en el mundo: una mirada actual y su futuro escenario epidemiológico:} El objetivo de este estudio fue evaluar el escenario epidemiológico contemporáneo del cáncer a nivel mundial y realizar una proyección del impacto de esta enfermedad en la población chilena para orientar la planificación sanitaria. Para ello, se realizó una revisión narrativa de la literatura fundamentada en datos oficiales de agencias internacionales (como la OMS, GLOBOCAN e IARC) e instituciones nacionales (como el MINSAL, DEIS y encuestas nacionales de salud), analizando el incremento histórico de mortalidad y morbilidad entre 1986 y 2018, mediante la recopilación de estadísticas de incidencia, mortalidad y factores de riesgo asociados a la oncología, con el fin de sintetizar indicadores precalculados provenientes de fuentes secundarias validadas. Como resultado, en 2018 se diagnosticaron 53365 nuevos casos de cáncer en Chile y la mortalidad por neoplasias aumentó un 109\% entre 1986 y 2016, proyectándose como una potencial causa de muerte en Chile. Se identificó que los cánceres más frecuentes fueron próstata, colorrectal, mama, estómago y pulmón, y se confirmó que, al momento del estudio (2020), el cáncer estaba desplazando a las enfermedades cardiovasculares como primera causa de muerte, siendo ya la principal causa en cinco regiones del país: Arica y Parinacota, Antofagasta, Coquimbo (mencionada como La Serena en la investigación), Los Lagos y Aysén. Por otro lado, las proyecciones estimaron un incremento del 77.6\% en nuevos diagnósticos para 2040 y se consideró que cerca del 40\% de los diagnósticos oncológicos están asociados a factores de riesgo modificables, concluyendo que es imperativo implementar un rediseño de las políticas de salud pública enfocado en el fomento de estilos de vida saludables, el control de la obesidad y la detección temprana de enfermedades o factores de riesgo. Cabe destacar que la validez del estudio se sustenta en la autoridad y confiabilidad de las fuentes citadas (\cite{scielo4}).
\end{itemize}

Cabe mencionar que, en la Tabla \ref{tab:resumen_scielo} se presentan los resumenes de cada uno de los artículos internacionales seleccionados sobre la caracterización de cohortes oncológicas mediante análisis epidemiológico y estadística descriptiva.

\begin{table}[htbp]
\centering
\footnotesize
\setlength{\tabcolsep}{2pt}
\renewcommand{\arraystretch}{1.15}
\caption{Resumen de estudios sobre caracterización de cohortes oncológicas mediante análisis epidemiológico y estadística descriptiva.}
\label{tab:resumen_scielo}
\begin{tabular}{|c|c|c|c|}
\hline
\parbox[t]{1.2cm}{\centering\textbf{Estudio}} & \parbox[t]{4.0cm}{\centering\textbf{Objetivo principal}} & \parbox[t]{5.0cm}{\centering\textbf{Datos y metodología}} & \parbox[t]{5.0cm}{\centering\textbf{Hallazgos principales}} \\ \hline

\parbox[t]{1.2cm}{\centering \cite{scielo1}} & 
\parbox[t]{4.0cm}{Caracterizar el perfil clínico y demográfico de pacientes con cáncer pulmonar y su asociación con factores de riesgo locales.} & 
\parbox[t]{5.0cm}{\textbf{Datos:} Fichas clínicas de 220 pacientes (Hospital Regional de Temuco). \newline \textbf{Métodos:} Estudio observacional y análisis de estadística descriptiva.} & 
\parbox[t]{5.0cm}{Predominio masculino y edad media de 68 años. Alta detección en etapas avanzadas (81.9\%). Asociación significativa entre expresión de PDL-1 y exposición a biomasa.} \\ \hline

\parbox[t]{1.2cm}{\centering \cite{scielo2}} & 
\parbox[t]{4.0cm}{Analizar la variación histórica y tendencias de mortalidad por cáncer oral y orofaríngeo en Chile (1955-2021).} & 
\parbox[t]{5.0cm}{\textbf{Datos:} Registros de defunción oficiales (INE, DEIS). \newline \textbf{Métodos:} Estudio descriptivo cuantitativo y análisis mediante regresión Joinpoint.} & 
\parbox[t]{5.0cm}{Clara tendencia al alza en la mortalidad. Mayor afectación en hombres, aunque la brecha de género disminuye. Evidencia la falta de políticas de detección temprana.} \\ \hline

\parbox[t]{1.2cm}{\centering \cite{scielo3}} & 
\parbox[t]{4.0cm}{Describir el comportamiento del cáncer de tiroides en pacientes sometidos a resolución quirúrgica.} & 
\parbox[t]{5.0cm}{\textbf{Datos:} Registros clínicos de 124 pacientes operados (Hospital Clínico Herminda Martín). \newline \textbf{Métodos:} Análisis descriptivo de variables clínicas y citología Bethesda.} & 
\parbox[t]{5.0cm}{Predominio en mujeres (86.2\%) y carcinoma papilar (93.1\%). El 36\% presentó invasión extratiroidea, variable asociada críticamente a los largos tiempos de espera.} \\ \hline

\parbox[t]{1.2cm}{\centering \cite{scielo4}} & 
\parbox[t]{4.0cm}{Evaluar el escenario epidemiológico del cáncer a nivel mundial y proyectar su impacto en la población chilena.} & 
\parbox[t]{5.0cm}{\textbf{Datos:} Estadísticas oficiales (OMS, GLOBOCAN, DEIS, MINSAL). \newline \textbf{Métodos:} Revisión narrativa y análisis de tasas de incidencia y mortalidad.} & 
\parbox[t]{5.0cm}{El cáncer desplazó a patologías cardiovasculares en cinco regiones. Se proyecta un aumento del 77.6\% en diagnósticos para 2040, urgiendo políticas de prevención.} \\ \hline

\end{tabular}
\end{table}


\subsection{Enfoques de solución}

Con base en la literatura (los estudios mencionados anteriormente), se identificaron tres enfoques analíticos aplicados al estudio de pacientes oncológicos. El primero, de \textbf{estadística tradicional}, emplea técnicas como la regresión logística y los modelos lineales, útiles para describir asociaciones lineales simples; sin embargo, asumen linealidad y aditividad entre las variables, lo que resulta insuficiente cuando los desenlaces en GRD dependen de interacciones no lineales y multidimensionales, patrones que la estadística clásica tiende a simplificar (\cite{article5}). 

El segundo enfoque, de \textbf{aprendizaje automático supervisado}, permite modelar relaciones complejas y descubrir patrones ocultos que la estadística tradicional omite, logrando clasificar desenlaces como severidad o mortalidad (\cite{article1}). No obstante, su principal limitación es la ausencia de una justificación clara de las predicciones, lo que dificulta su interpretación y validación clínica, siendo incompatible con el objetivo de apoyar la toma de decisiones estratégicas.

El tercer enfoque, de \textbf{aprendizaje automático explicable (XAI)}, combina la precisión predictiva de los modelos de aprendizaje automático con la capacidad interpretativa, alineándose con la necesidad de generar perfiles y evidencia sanitaria accionable para instituciones como FONASA, en donde la explicabilidad se entiende como la cuantificación de la contribución marginal de cada variable a la predicción final, utilizando SHAP (SHapley Additive exPlanations) y midiendo su aporte en la identificación de reglas de decisión contrastadas con estudios científicos, con el fin de evaluar su coherencia clínica.


\subsection{Justificación del enfoque seleccionado}

Para esta propuesta, se seleccionó el enfoque de \textbf{aprendizaje automático supervisado explicable (XAI)} por su capacidad única para abordar la problemática central de este estudio, al permitir construir perfiles operativos que identifican qué combinaciones de factores clínicos, demográficos y hospitalarios determinan los distintos niveles de mortalidad, severidad y consumo de recursos en episodios oncológicos, transformando un problema técnico de clasificación en una herramienta de inteligencia sanitaria. 

Además, aunque los datos GRD no fueron diseñados para la investigación clínica, este enfoque resulta pertinente porque el aprendizaje automático permite modelar y descubrir relaciones complejas entre las variables de entrada y las de salida, las cuales tienden a ser simplificadas por los métodos estadísticos tradicionales. Adicionalmente, la capa de explicabilidad posibilita validar la coherencia clínica de los patrones detectados o revelar hallazgos novedosos, mientras que el carácter supervisado se justifica en la definición de las variables de salida de interés en el conjunto de datos.

Por último, este enfoque aporta insumos prácticos para la gestión sanitaria, ya que la interpretabilidad permite implementar estrategias basadas en la identificación de factores específicos que elevan cada riesgo, apoyando la planificación sustentada en evidencia. De esta manera, se avanza desde asignaciones basadas en promedios históricos hacia una planificación ajustada a la complejidad real del sistema público, optimizando la respuesta de la red frente a la demanda oncológica.

